\refstepcounter{section}
\section*{Lecture 19: Trace and Contraction Identities}

\addcontentsline{toc}{section}{Lecture 19: Trace and Contraction Identities}
In the previous section, we saw that calculating the trace of the Dirac matrices become really important since these define the inner products. Let us look at some of the trace identities followed by the Dirac matrices.
\subsection{Trace Identitites}
\begin{itemize}
    \item $\Tr(\gamma^\mu)=0$\\[0.2cm]
    \textit{Proof.} We defined $\gamma^0=\beta$ which was already traceless. From the anti-commutation relation, we know that $\beta \alpha_i = -\alpha_i \beta$ and taking trace on both sides, we get $\Tr \gamma^i = 0$
    \item $\{\gamma^\mu,\gamma^\nu\}=2\eta^{\mu\nu}\mathds{1}$\\[0.2cm]
    \textit{Proof.} Let us take the spatial indices first. From the definition of $\alpha_i$ and $\beta$, we have the following:
    \begin{align*}
        \gamma^i\gamma^j+\gamma^j\gamma^i = \beta (\alpha_i \beta) \alpha_j + i\leftrightarrow j=-\beta  (\beta \alpha_i)\alpha_j+ i\leftrightarrow j=-(\alpha_i\alpha_j +\alpha_j\alpha_i ) =-2\delta_{ij}\mathds{1}=2\eta^{ij}\mathds{1}
    \end{align*}
    Now, consider one temporal index, that is,
    \begin{align*}
        \{\gamma^0, \gamma^i\}= \beta^2 \alpha_i +\beta \alpha_i\beta =\alpha_i -\alpha_i\beta^2 =0
    \end{align*}
    Finally, we have $\{\gamma^0, \gamma^0\} = 2\beta^2 = 2\mathds{1}=2\eta^{00}\mathds{1}$. Combining the above, we have the identity. \\[0.2cm]
    If we instead take the anti-commutation of the covariant indices, then also the identity remains similar, that is,
    $$\{\gamma_{\mu}, \gamma_{\nu}\} = 2\eta_{\mu\nu}\mathds{1}$$
    \item $\Tr(\gamma^\mu\gamma^\nu)=4\eta^{\mu\nu}$\\[0.2cm]
    \textit{Proof.} From the anti-commutation relation of the Dirac matrices, we have:
    \begin{align*}
      2\Tr(\gamma^\mu\gamma^\nu) =  \Tr\{\gamma^\mu,\gamma^\nu\} = \Tr (2\eta^{\mu\nu}\mathds{1}) =2\eta^{\mu\nu} \times 4 \implies \Tr(\gamma^\mu\gamma^\nu) = 4 \eta^{\mu\nu}
    \end{align*}
    \item $(\gamma^5)^2 = \mathds{1}$\\[0.2cm]
    \textit{Proof.} From the definition of $\gamma^5$ and the anti-commutation between the gamma matrices, we have,
    \begin{align*}
        \gamma^5 \gamma^5 &= - (\gamma^0 \gamma^1 \gamma^2 \gamma^3 )(\gamma^0 \gamma^1 \gamma^2 \gamma^3 )\\
        &=\gamma^1\gamma^2\gamma^3(\gamma^0)^2 \gamma^1 \gamma^2\gamma^3\\ 
        &=\gamma^1\gamma^2\gamma^3\gamma^1 \gamma^2\gamma^3\\
        &=\gamma^2\gamma^3 (\gamma^1)^2\gamma^2\gamma^3\\
        &=-\gamma^2\gamma^3 \gamma^2 \gamma^3 \\
        &= \gamma^3( \gamma^2)^2 \gamma^3 \\
        &= -(\gamma^3)^2\\
        &= \mathds{1}
    \end{align*}
    \item $\gamma^5$ is Hermitian. \\[0.25cm]
    \textit{Proof.} We will use the fact that $\qty(\gamma^i)^\dagger = -\gamma^i$ for $i=1,2,3$ and $\qty(\gamma^0)^\dagger = \gamma^0$ along with the anti-commutation relation. \begin{align*}
        (\gamma^5)^\dagger = -i \qty(\gamma^3)^\dagger \qty(\gamma^2)^\dagger \qty(\gamma^1)^\dagger \qty(\gamma^0)^\dagger = i \gamma^3 \gamma^2 \gamma^1 \gamma^0 
    \end{align*}
    To bring $\gamma^0$ to correct position (in front), we need 3 flips, for $\gamma^1$ we need 2 flips and for $\gamma^2$ we need only one flip, hence total $3+2+1=6$ flips. Each flip introduces a negative sign, so total $(-1)^6=1$, hence $(\gamma^5)^\dagger = i\gamma^0 \gamma^1 \gamma^2 \gamma^3 =\gamma^5$.
    \item $\{\gamma^5, \gamma^\mu\} = 0 $\\[0.25cm]
    \textit{Proof.} Note that, in $\gamma^5\gamma^\mu = i \gamma^0 \gamma^1 \gamma^2 \gamma^3 \gamma^\mu$, for exchanging each $\mu\neq \nu$, we introduce a $-1$. Now, since there are $4$ indices, $\mu$ must be one of these four indices and hence, we have the same index together, we do not need the $-1$. Thus, we introduce exactly $(-1)^3 = -1$ while transporting $\gamma^\mu$ to the front. Then we have $\gamma^5\gamma^\mu = -\gamma^\mu\gamma^5$, thus proving the identity. 
    \item $\Tr (\gamma^{\mu_1 }\gamma^{\mu_2 }\ldots \gamma^{\mu_{2n+1} })=0$ \\[0.25cm]
    \textit{Proof.} As $(\gamma^5)^2 = \mathds{1}$, let us multiply $\gamma^5$ to the above expression,
    \begin{align*}
        \gamma^{\mu_1 }\gamma^{\mu_2 }\ldots \gamma^{\mu_{2n+1} } &= \gamma^5\gamma^5\gamma^{\mu_1 }\gamma^{\mu_2 }\ldots \gamma^{\mu_{2n+1} }\\
    \end{align*}
    Shifting one of the $\gamma^5$ to the far right across the $2n+1$ indices will introduce $(-1)^{2n+1}$ and thus, we get:
    $$\gamma^{\mu_1 }\gamma^{\mu_2 }\ldots \gamma^{\mu_{2n+1} } = (-1)^{2n+1}\gamma^5\gamma^{\mu_1 }\gamma^{\mu_2 }\ldots \gamma^{\mu_{2n+1} }\gamma^5$$
    Taking trace both sides and taking the cyclic property of the trace, we obtain the following. Thus, the product of odd number of gamma matrices is zero. 
    \item $\Tr(\gamma^\kappa \gamma^\lambda \gamma^\mu \gamma^\nu) = 4\eta^{\kappa\lambda}\eta^{\mu\nu}-4\eta^{\kappa\mu}\eta^{\lambda\nu}+4\eta^{\kappa\nu}\eta^{\lambda\mu}$\\[0.25cm]
    \textit{Proof.} Using the anti-commutation relation, we note the following: 
    \begin{align*}
        \{\gamma^\kappa, \gamma^\lambda\gamma^\mu\gamma^\nu\} &= \gamma^\kappa\gamma^\lambda\gamma^\mu\gamma^\nu +\gamma^\lambda\gamma^\mu\gamma^\nu\gamma^\kappa \\
        &=\gamma^\kappa\gamma^\lambda\gamma^\mu\gamma^\nu  + \mathcolor{red}{\gamma^{\lambda}\gamma^{\kappa}\gamma^{\mu}\gamma^{\nu}}-\mathcolor{red}{\gamma^{\lambda}\gamma^{\kappa}\gamma^{\mu}\gamma^{\nu}}+ \gamma^\lambda\gamma^\mu\gamma^\nu\gamma^\kappa \\
        &=\{\gamma^\kappa, \gamma^\lambda\}\gamma^\mu\gamma^\nu - \gamma^{\lambda}\gamma^{\kappa}\gamma^{\mu}\gamma^{\nu} -\mathcolor{blue}{\gamma^{\lambda}\gamma^{\mu}\gamma^{\kappa}\gamma^{\nu}} +\mathcolor{blue}{\gamma^{\lambda}\gamma^{\mu}\gamma^{\kappa}\gamma^{\nu}} + \gamma^\lambda\gamma^\mu\gamma^\nu\gamma^\kappa \\
        &=\{\gamma^\kappa, \gamma^\lambda\}\gamma^\mu\gamma^\nu - \gamma^\lambda\{\gamma^\kappa, \gamma^\mu\}\gamma^\nu +\gamma^\lambda\gamma^\mu\{\gamma^\kappa,\gamma^\nu\}\\
        &=2\eta^{\kappa\lambda}\gamma^\mu\gamma^\nu - 2\eta^{\kappa\mu}\gamma^\lambda\gamma^\nu +2\eta^{\kappa\nu}\gamma^\lambda\gamma^\mu
    \end{align*}
    Now, note that the required quantity is basically half the trace of the above anti-commutator.
    \begin{align*}
     \Tr({{\gamma^\kappa} {\gamma^\lambda}  \gamma^\mu  \gamma^\nu}) =  \frac{1}{2} \Tr \{\gamma^\kappa, \gamma^\lambda\gamma^\mu\gamma^\nu\}&=\Tr \qty(\eta^{\kappa\lambda}\gamma^\mu\gamma^\nu - \eta^{\kappa\mu}\gamma^\lambda\gamma^\nu +\eta^{\kappa\nu}\gamma^\lambda\gamma^\mu)\\
&=\eta^{\kappa\lambda}\Tr \gamma^\mu\gamma^\nu -\eta^{\kappa\mu} \Tr \gamma^\lambda\gamma^\nu +\eta^{\kappa\nu}\Tr \gamma^\lambda\gamma^\mu\\
&= 4 \eta^{\kappa\lambda}\eta^{\mu\nu}  -4\eta^{\kappa\mu}\eta^{\lambda\nu} + 4\eta^{\kappa\nu}\eta^{\lambda\mu}
    \end{align*} 
\item $\Tr \gamma^5 =0$\\[0.25cm]
\textit{Proof.} Since the Minkowski metric is entire diagonal, $\eta^{ij}=0$ for $i\neq j$ from which the identity follows using the previous identity. 
\item $\Tr(\gamma^5 \gamma^\mu) =0$\\[0.25cm]
\textit{Proof.} Note that $\mu\in \{0,1,2,3\}$ and $\gamma^5$ contains of all these indices. So, we can bring $\gamma^\mu$ adjacent to the same index in $\gamma^5$ by some flips and then it will give $(\gamma^\mu)^2$ which gives $\pm 1$ atmost and the remaining will be a product of odd number of gamma matrices whose trace we had found out to be zero. 
\item $\Tr(\gamma^5 \gamma^\mu \gamma^\nu) =0$\\[0.25cm] 
\textit{Proof.} We will use the fact the $\gamma^5$ anti-commutes with every $\gamma^\mu$. Let us introduce two factors of $\gamma^\alpha$ in the trace such that $\alpha\neq \mu,\nu$. Note that $(\gamma^\alpha)^2$ will be $1$ if $\alpha=0$, else, it will be $-1$. In any case, the trace will change by $\pm 1$ 
\begin{align*}
    \Tr(\gamma^5 \gamma^\mu \gamma^\nu) = (-1)^{1-\delta_{\alpha, 0}}\Tr(\gamma^5 \gamma^\mu \gamma^\nu\gamma^\alpha \gamma^\alpha) &=(-1)^{1-\delta_{\alpha, 0}} (-1)^3\Tr(\gamma^\alpha \gamma^5  \gamma^\mu \gamma^\nu \gamma^\alpha) \\ &= -(-1)^{1-\delta_{\alpha, 0}} \Tr( \gamma^5  \gamma^\mu \gamma^\nu \gamma^\alpha\gamma^\alpha) \\ &= -\Tr(\gamma^5 \gamma^\mu \gamma^\nu)\\ \implies \Tr(\gamma^5 \gamma^\mu \gamma^\nu)=0
\end{align*}
In the last step, we had used the cyclic property of the trace to bring $\gamma^\alpha$ to the back again. 
\item $\gamma^5 \gamma^\mu\gamma^\nu \gamma^\alpha \gamma^\beta = -4i \epsilon^{\mu\nu\alpha\beta}$\\[0.25cm]
\textit{Proof.} If any two of $\mu,\nu,\alpha,\beta$ are same, we could bring them together by flipping, introducing atmost $\pm 1$. The remaining will be the product of $\gamma^5$ and two other gamma matrices, whose trace we had already shown to be zero.\\[0.2cm] Thus, let us assume that all the indices are different. Then in the product  $\gamma^5 \gamma^\mu\gamma^\nu \gamma^\alpha \gamma^\beta$, each gamma will appear twice. Hence we flip them such that all same index gamma appear together which would give us one of $\pm \mathds{1}$, hence the trace will be of the form $\pm i \Tr \mathds{1} = \pm 4i$. \\[0.25cm]
Suppose $\mu,\nu,\alpha,\beta$ appear as an even permutation of $0,1,2,3$. Then we have, $$\Tr (\gamma^5\gamma^\mu\gamma^\nu \gamma^\alpha \gamma^\beta) =-i \Tr(\gamma^5\times \mathcolor{red}{i\gamma^\mu\gamma^\nu \gamma^\alpha \gamma^\beta})$$ Now the red factor can be brought to the form $\gamma^5$ using an even number of exchanges. Since we saw that $(\gamma^5)^2=\mathds{1}$, we have $\Tr (\gamma^5\gamma^\mu\gamma^\nu \gamma^\alpha \gamma^\beta) = -4i$. \\[0.2cm]
If $\mu,\nu,\alpha,\beta$ appear as an odd permutation of $0,1,2,3$, then the same logic follows, however, the red factor introduces a negative sign since in this case, by an odd number of exchanges, the factor can be brought to the form of $\gamma^5$. Hence the trace will give $4i$. \\[0.2cm]
As a summary, we saw that if any two of the four indices are equal, the trace is zero. If indices are even permutation of $0,1,2,3$ then trace $\to\ -4i$ else if odd permutation, then trace $\to\ 4i$. This is essentially the same statement as the contravariant Levi-Civita tensor defined as, 
$$\epsilon^{\alpha\beta\kappa\lambda } = \begin{cases}
    +1 & \alpha\beta\kappa\lambda \to \ \text{even-permutation of } 0,1,2,3\\
    -1 &  \alpha\beta\kappa\lambda \to \ \text{odd-permutation of } 0,1,2,3\\
    0 & \text{else}
\end{cases}$$
Thus, we prove the identity. 
\end{itemize}
\subsection{Contraction Identities}
Let us now look at some contraction identities where the Dirac matrices will be contracted. Unlike the trace identitities which gave a number as a final answer, these will give us some matrices. 
\begin{itemize}
    \item $\gamma^\mu \gamma_\mu = 4\mathds{1}$\\[0.25cm]
    \textit{Proof.} Since Dirac matrices can be treated as four vectors, we can lower and raise their indices with the metric $\eta$. Then, 
    \begin{align*}
        \gamma^\mu \gamma_\mu &= \eta_{\mu\nu}\gamma^{\nu}\gamma^{\mu}\\
        &=\frac{1}{2}\qty(\eta_{\mu\nu} + \eta_{\nu\mu})\gamma^{\nu}\gamma^{\mu} \qquad (\text{as metric is symmetric})\\
        &=\frac{1}{2}\qty(\eta_{\mu\nu}\gamma^{\nu}\gamma^{\mu}  + \eta_{\nu\mu}\gamma^{\nu}\gamma^{\mu} )\\
        &=\frac{1}{2}\qty(\eta_{\mu\nu}\gamma^{\nu}\gamma^{\mu}  + \eta_{\mu\nu}\gamma^{\mu}\gamma^{\nu} )\qquad (\text{relabelling dummy indices})\\ 
        &=\frac{1}{2}\eta_{\mu\nu} \qty(\gamma^{\nu}\gamma^{\mu}  + \gamma^{\mu}\gamma^{\nu} )\\
        &=\frac{1}{2}\eta_{\mu\nu}  \times 2\eta^{\mu\nu}\mathds{1}\qquad (\text{from anti-commutation of gamma matrices})\\
        &=\tensor{\delta}{^\mu _\mu}\mathds{1}\\
        &=4\mathds{1}
    \end{align*}
    \item $\gamma^\mu \gamma^\nu \gamma_\mu = -2\gamma^\nu$\\[0.25cm]
    \textit{Proof.} Using the anti-commutation relation and the previous identity, we will derive this.
    \begin{align*}
        \gamma^\mu \gamma^\nu \gamma_\mu &=\gamma^\mu \gamma^\nu \eta_{\mu\beta}\gamma^\beta\\
        &=\gamma^\mu \eta_{\mu\beta}\qty(2\eta^{\nu\beta}-\gamma^\beta \gamma^\nu)\\
        &= 2\gamma^\mu \eta_{\mu\beta}\eta^{\nu\beta} - \gamma^\mu \gamma^\beta \eta_{\mu\beta}\gamma^\nu\\
        &=2\gamma^\mu\tensor{\delta}{^\nu _\mu} - \underbrace{\gamma^\mu \gamma_\mu}_{4\mathds{1}} \gamma^\nu\\
        &=2\gamma^\nu - 4\gamma^\nu\\
        &=-2\gamma^\nu
    \end{align*}
    \item $\gamma^\mu \gamma_\nu \gamma_\lambda \gamma_\mu = 4\eta_{\lambda \nu}\mathds{1}$\\[0.25cm]
    \textit{Proof.} We will use the anti-commutation relation and the previous identity together to reduce the expression. 
    \begin{align*}
        \gamma^\mu \gamma_\nu \gamma_\lambda \gamma_\mu &= \gamma^\mu \gamma_\nu \qty(2\eta_{\lambda\mu} - \gamma_\mu \gamma_\lambda)\\
        &=2\eta_{\lambda\mu}\gamma^{\mu}\gamma_{\nu} - \gamma^\mu \gamma_\nu \gamma_\mu \gamma_\lambda\\
        &=2\gamma_\lambda \gamma_\nu - \eta_{\alpha\nu}\underbrace{\gamma^{\mu}\gamma^{\alpha}\gamma_\mu}_{-2\gamma^\alpha} \gamma_\lambda\\
        &=2\qty(\gamma_\lambda \gamma_\nu + \gamma_\nu\gamma_\lambda)\\
        &=4\eta_{\lambda\nu}\mathds{1}
    \end{align*}
    \item $\gamma^\mu \gamma_\nu \gamma_\lambda \gamma_\rho \gamma_\mu = -2 \gamma_\rho \gamma_\lambda \gamma_\nu$\\[0.25cm]
    \textit{Proof.} We will again use the anti-commutation relation and the previous identity for this. 
    \begin{align*}
        \gamma^\mu \gamma_\nu \gamma_\lambda \gamma_\rho \gamma_\mu &= \gamma^\mu \gamma_\nu \gamma_\lambda \qty(2\eta_{\rho\mu}-\gamma_\mu \gamma_\rho) \\
        &= 2 \gamma_\rho \gamma_\nu \gamma_\lambda - \underbrace{\gamma^\mu \gamma_\nu \gamma_\lambda \gamma_\mu}_{4\eta_{\lambda\nu}\mathds{1}} \gamma_\rho\\ 
        &=2\gamma_\rho\qty(\gamma_\nu \gamma_\lambda - 2\eta_{\lambda\nu})\\
        &= -2\gamma_\rho \times ( 2\eta_{\lambda\nu} -\gamma_\nu \gamma_\lambda) \\
        &= -2\gamma_\rho \gamma_\lambda \gamma_\nu
    \end{align*}
\end{itemize}
These are few of the contraction and trace identities that we checked. There exists many such other identities and we can just play with the Dirac matrices to find them \emoji{wink}!